% =====================================================================
% Self-contained section fragment: the McCormick blow-up and the
% VM-One-Armed-Rescue solver.  Ready to \input into neurips_2026.tex.
% Relies on preamble macros \Tmax \Tmin \CGM \Vone \ie \eg and on
% cleveref/natbib; cites keys already present in ref.bib.  Any genuinely
% new keys would live in update/newrefs_solver.bib.
% =====================================================================
\section{The McCormick Blow-Up and the VM-One-Armed-Rescue Solver}
\label{sec:solver_vm}

The expressive power of the joint program is not free. Turning the tiling
choice into a decision variable (\cref{sec:v1}) makes the DRAM-savings
accounting \emph{bilinear}; linearizing it exactly through McCormick
envelopes~\citep{mccormick1976computability} then multiplies the size of the
constraint matrix and drives its conditioning across roughly ten orders of
magnitude. We describe this blow-up (\cref{sec:mccormick_blowup}) and the
custom first-order GPU solver that renders the resulting linear program
tractable at LLM scale (\cref{sec:vm_solver}). The two are inseparable: we
adopt the exact linearization \emph{only because} we have a solver engineered
to consume its output.

\subsection{The McCormick Blow-Up}
\label{sec:mccormick_blowup}

\paragraph{A bilinear coupling.}
In the sequential pipeline the per-configuration timing $\Tmax_i(c)$,
$\Tmin_i(c)$ and per-tensor DRAM cost $t_i^k(c)$ are fixed constants. \Vone{}
instead selects the tiling inside the LP through one-hot variables
$y_{i,c}\in[0,1]$ with $\sum_{c\in\mathcal{C}_i} y_{i,c}=1$
(\cref{sec:v1}). The savings realized when tensor $k$ of operation $i$ is
pinned on chip are only credited if that tensor is \emph{both} selected under
configuration $c$ \emph{and} resident, so the DRAM-savings term becomes a
product of two decision variables,
\begin{equation}
\label{eq:bilinear_coupling}
t_i^k(c)\; \underbrace{y_{i,c}}_{\text{tiling selection}}\cdot
\underbrace{p_{\text{assoc}(i,k),\,o(i)}}_{\text{residency}},
\end{equation}
where $p_{e,t}\in[0,1]$ is the edge-residency variable and $o(i)$ the
execution order of $i$. The bilinear factor $y_{i,c}\,p$ is nonconvex and
cannot be placed in a linear program directly.

\paragraph{The McCormick box.}
For each such product we introduce an auxiliary variable
$w_{i,c,k} := y_{i,c}\, p_{\text{assoc}(i,k),\,o(i)}$ and bound it inside a box
cut out by four linear inequalities~\citep{mccormick1976computability},
\begin{equation}
\label{eq:mcc_box}
w_{i,c,k} \le y_{i,c},\qquad
w_{i,c,k} \le p_{\text{assoc}(i,k),\,o(i)},\qquad
w_{i,c,k} \ge y_{i,c} + p_{\text{assoc}(i,k),\,o(i)} - 1,\qquad
w_{i,c,k} \ge 0 .
\end{equation}
Geometrically, the four planes \emph{bracket} the bilinear surface $w=y\,p$
over the unit box $[0,1]^2$: $w\le y$ and $w\le p$ form an over-estimating
upper envelope, $w\ge 0$ and $w\ge y+p-1$ an under-estimating lower envelope,
and the feasible $w$ is squeezed between them. Because both factors lie in
$[0,1]$, this box is the \emph{tight convex hull} of the graph of $y\,p$: the
smallest convex relaxation over the interior, and at each of the four integer
vertices $(y,p)\in\{0,1\}^2$ the upper and lower planes pinch shut so that
$w_{i,c,k}=y_{i,c}\,p$ holds \emph{exactly}, with no relaxation gap. Since
$y_{i,c}$ is one-hot and residency rounds to $\{0,1\}$, the solution lands
only on these gap-free corners, so \eqref{eq:mcc_box} introduces no tuning
constants and no bias. The same construction handles the two tile-to-residency
families $z^{\mathrm{act}}_{e,c,t}$, $z^{W}_{i,c,t}$ of the tile-dependent
LP \eqref{eq:v1}.

\paragraph{A matrix that explodes.}
Exactness is bought with size. Each bilinear product spends one auxiliary
variable and four inequalities (three coupling rows plus a nonnegativity
bound), replicated over every $(\text{operation},\text{tile
config},\text{tensor})$ triple. Summed over the network this dominates the
program. For DeepSeek-V3 the \Vone{} LP reaches roughly $2{,}504{,}019$
variables, $2{,}983{,}450$ constraints, and $6{,}094{,}156$ nonzeros, of which
the McCormick auxiliaries and their selection partners are the large majority.
Worse than its size, the matrix is severely \emph{ill-conditioned}: the
objective and constraint coefficients span about ten orders of magnitude,
because Timeloop per-tile cycle counts run from $10^{3}$ to $10^{10}$ while the
residency and byte terms are $O(1)$ and the roofline memory term scales as
$1/\text{BW}$; the products injected by \eqref{eq:mcc_box} sit at the extreme
ends of this range simultaneously. State it plainly: this blow-up is
acceptable \emph{only} because we have a solver that can handle a
multi-million-variable, ten-orders-ill-conditioned LP in seconds. Standard
tools cannot. The remainder of this section is that solver.

\subsection{The Variable-Metric-One-Armed-Rescue Solver}
\label{sec:vm_solver}

We solve the McCormick-blown-up LP with a custom, matrix-free, first-order GPU
solver written in JAX, which we call \textbf{VM-One-Armed-Rescue}. Its
backbone is restarted \emph{primal-dual hybrid gradient}
(PDHG)~\citep{chambolle2011first} in the practical PDLP
lineage~\citep{applegate2025pdlp,applegate2026pdlp}: every iteration is a pair
of sparse matrix-vector products and elementwise projections costing
$\mathcal{O}(\mathrm{nnz})$, with no factorization ever formed. This is what
lets a two-million-variable LP live entirely on a single GPU. Three
ingredients turn the generic restarted PDHG of GPU solvers such as
cuPDLP~\citep{lu2024cupdlpjlgpuimplementationrestarted,%
lu2025cupdlpxenhancedgpubasedfirstorder} into a method that actually converges
on our instances:
\begin{itemize}[leftmargin=1.5em, itemsep=1pt, topsep=2pt]
  \item \textbf{Diagonal variable-metric rescaling.} Rather than a fixed
  Euclidean metric, we precondition primal and dual iterates with a diagonal
  variable metric, initialized by Ruiz-style
  equilibration~\citep{ruiz2001scaling} of the coefficient matrix and then
  driven online by a smoothed estimate of
  $\sqrt{r_{\mathrm{dual}}/r_{\mathrm{primal}}}$, the square-root ratio of dual
  to primal residual. This rebalances the two sides of the saddle point that
  the ten-orders coefficient spread would otherwise pull apart.
  \item \textbf{One-sided ``rescue'' restart.} Stock PDLP restarts the primal
  and dual symmetrically; on our ill-conditioned LPs that heuristic overshoots
  and the duality gap oscillates without settling. VM-One-Armed-Rescue instead
  restarts \emph{asymmetrically}, re-anchoring only the side (the ``one arm'')
  whose residual is lagging, so the trailing side is rescued without
  destabilizing the side that is already converging.
  \item \textbf{Halpern anchoring.} Restarts use Halpern-style
  anchoring~\citep{halpern1967fixed} toward the last restart point, which
  contracts the iterates and yields the stable, monotone convergence that the
  bare restart heuristic lacks.
\end{itemize}
A JIT-compiled \texttt{lax.scan} batches ${\sim}10^{5}$ inner iterations onto
the accelerator at a few microseconds each and fuses the equilibration loop
into a single kernel.

\paragraph{Why it matters.}
On the \Vone{}-scale co-design LP, off-the-shelf solvers do not deliver a
usable solution: CPU simplex and barrier (HiGHS) exhaust factorization memory
or stall on the badly scaled matrix, and stock first-order GPU solvers
(cuPDLP) leave the duality gap oscillating rather than closing it.
VM-One-Armed-Rescue solves the same instances to primal--dual convergence in
\emph{seconds} on one GPU (the eight production blueprints report the
custom JAX PDHG solving the \Vone{} LP in $1.7$--$6.7$\,s), landing within a
bounded, reported rounding gap of the mature HiGHS optimum wherever HiGHS can
reach one. To our knowledge this is the first solver to make the
McCormick-linearized \emph{joint} tiling-plus-fusion-plus-placement co-design
LP tractable at LLM scale, and that tractability is precisely what makes the
inner LP cheap enough to sit inside a thousands-of-trials outer search rather
than a one-off solve.

\paragraph{Scope of the claim.}
We state the solver claim exactly where it holds. VM-One-Armed-Rescue is
engineered for, and validated on, the \Vone{} and tile-dependent LLM-scale LPs
of the ${\sim}2.5$M-variable class characterized above (the LPs that emit
the blueprints), and it is on those instances that the contrast with HiGHS
above is measured. We do \emph{not} claim it solves the strictly larger
monolith that additionally lifts the chip choice into a single program: the
Tier-1 hardware-menu LP for DeepSeek-V3 reaches roughly $20.7$M variables and
$33.3$M constraints, and on that instance VM-One-Armed-Rescue did not converge;
neither did HiGHS, which stalled at the time limit without reaching
primal feasibility. That both a GPU first-order method and a mature
simplex/barrier code fail on the menu monolith is exactly what motivates
decomposing it rather than solving it whole; the ${\sim}2.5$M-variable
per-workload LP is the tractable unit the outer search is built on.
